
\subsection{Where Initiation Fits}
\label{sec:placement}

Physical safety---collisions, forces, joint limits---remains
necessary
\citep{akalin2021,zhang2025safevla,kojima2025fmr}, and cyber and
privacy safety---data leakage, jailbreak prompting---are also relevant
for foundation-model robots
\citep{robey2024jailbreaking,bai2022constitutional}. Beyond these,
robots that initiate contact in public or domestic spaces face a third
question: whether to act at all.

HRI already distinguishes physical from perceived safety
\citep{akalin2021,weiss2023perceived}, and social-error taxonomies
list what goes wrong once interaction is already underway
\citep{tian2021socialerrors}. What is missing is a layer for the
earlier decision of whether the robot should start in the first place
\citep{shi2011spatial,edirisinghe2024shopworker}. We call this layer
\emph{initiation authorization} and place it as a critical subclass of
social-interaction safety.

\begin{table}[H]
\footnotesize
\centering
\caption{Safety classes for generalist robots.}
\label{tab:taxonomy}
\begin{tabular*}{\columnwidth}{@{\extracolsep{\fill}}>{\raggedright\arraybackslash}p{0.21\columnwidth}>{\raggedright\arraybackslash}p{0.37\columnwidth}>{\raggedright\arraybackslash}p{0.26\columnwidth}@{}}
\toprule
\textbf{Class} & \textbf{Safety focus} & \textbf{Example failure} \\
\midrule
Physical & Contact force, clearance & Unintended touch \\
\midrule
Cyber / privacy & Access control, PII & Leaked session data \\
\midrule
Social-interaction & Felt safety, norms & Interrupt, pressure \\
\midrule
\quad$\hookrightarrow$ Initiation & When to act, whom to address & Premature commitment \\
\bottomrule
\end{tabular*}
\end{table}

Initiation also spans more than speech. On talking platforms the first
commitment is often a word. On manipulation-centric VLAs it can be
the first grasp of a contested object, an uninvited entry into
personal space, or a clarifying question that cuts across a human
conversation \citep{rt2-2023,openvla-2024,huang2016anticipatory}.
What matters is \emph{irreversibility under uncertainty, not the
modality}. Generalist safety also depends on context, intent, and
human expectation, not on constraint enforcement alone.
Table~\ref{tab:themes} maps these themes to initiation decisions.

\begin{table}[H]
\small
\centering
\caption{Safety themes mapped to pre-commitment initiation.}
\label{tab:themes}
\begin{tabular*}{\columnwidth}{@{\extracolsep{\fill}}>{\raggedright\arraybackslash}p{0.32\columnwidth}>{\raggedright\arraybackslash}p{0.62\columnwidth}@{}}
\toprule
\textbf{Theme} & \textbf{Initiation analogue} \\
\midrule
Context (egg/knife) & Same \textsc{speak}/reach OK in demo, intrusive at home \\
\midrule
Intent & ``Looking at robot'' $\neq$ willing to be addressed \\
\midrule
Preference & Bold greeting vs.\ privacy-seeking visitor \\
\midrule
Asymmetric harm & Awkward interrupt vs.\ missed visitor \\
\midrule
FM / data risk & Who may speak, act, log, or adapt from traces \\
\bottomrule
\end{tabular*}
\end{table}

\FloatBarrier

\subsection{From Specialist to Generalist Safety}

Specialist deployments embed assumptions that generalists cannot
inherit: (i) one threshold $\theta$ works for everyone,
(ii) an engagement score reveals intent without any probe response
\citep{salam2017},
(iii) speaking too soon and waiting too long are symmetric mistakes, and
(iv) the same script applies across all venues. Generalists need
\emph{parametric authorization}: a per-venue and per-user dial on risk
posture that does not require retraining the whole VLA stack
\citep{liu2025embodied}. Table~\ref{tab:specialist} contrasts the two
regimes.

\begin{table}[H]
\small
\centering
\caption{Specialist vs.\ generalist safety assumptions.}
\label{tab:specialist}
\begin{tabular*}{\columnwidth}{@{\extracolsep{\fill}}lll@{}}
\toprule
 & \textbf{Specialist} & \textbf{Generalist} \\
\midrule
Harm model & Workspace, force & Norms, timing \\
\midrule
Observable state & Geometry & Latent receptivity \\
\midrule
Safe action & Stay in envelope & Sometimes refrain \\
\midrule
Tuning & Calibrate once & Per-venue dial \\
\bottomrule
\end{tabular*}
\end{table}

\subsection{Layered Safety and Adjacent Literatures}

If specialist assumptions fail, a single safety layer is also
insufficient. A unified picture has three shields, not one.
Physical shields and shielding-aware motion planning
\citep{hu2022sharp} handle contact risk.
VLA alignment \citep{zhang2025safevla,kojima2025fmr} constrains
plans once acting is allowed.
An \emph{initiation shield} \citep{ravichandran2025roboguard} sits
before the first irreversible social commitment.
Recent FMR guardrails reduce unsafe plans and resist jailbreaks, but
they assume the robot is already permitted to attempt the commanded
action. They do not decide whether the social channel should be opened
\citep{kojima2025fmr,ravichandran2025roboguard,robey2024jailbreaking}.
We now distinguish initiation authorization from three adjacent
literatures.

\paragraph{vs.\ physical safety and VLA guardrails.} Collision
avoidance and SafeVLA-style constraints ask: will the motion harm
someone? \citep{zhang2025safevla} Initiation authorization asks:
should the robot make an irreversible social commitment now? A
collision-free robot can still be socially harmful. A near-miss
with a knife is the physical analogue
\citep{akalin2021}.

\paragraph{vs.\ preference alignment and LLM safety.} Preference
alignment shapes what a model says or does after a channel is open
\citep{bai2022constitutional}. Embodied cold-start differs:
low-cost probes change the human's state.
Harm is social and bidirectional, and the irreversible commitment
precedes any human correction. Red-teaming a dialogue model does not
substitute for probe--response dynamics on a physical robot
\citep{robey2024jailbreaking}.

\paragraph{vs.\ engagement and proactive HRI.} Proactive-robot
research studies \emph{how} to initiate well---timing, position,
naturalness
\citep{broek2024proactive,hoffman2007anticipatory,shi2011spatial,bergstrom2008encounters}.
Engagement metrics maximize interaction quality
\citep{rich2010,salam2017}. We study \emph{whether} initiating is
permitted under uncertainty. Engagement scores do not constitute
consent. A safety policy must allow refraining when evidence is thin.
An engagement-optimal policy may act early. An initiation-safe one may
not.
